%%=====================================================================
%% The Icosian Completion of K_8
%%   Sol Monodromy Forces the Hurwitz and Trace-phi Enlargements:
%%   E_8 as a Derived Arithmetic Lattice
%%
%% Standalone derivation paper within the Master Equation Framework
%% corpus (outside the numbered series, per PI decision June 2026).
%% Version: v2 (June 2026) -- canonical mef_macros.tex adopted; \OrderMEF naming;
%%   dilaton notation corrected to \QCD per macro set
%%
%% Provenance:
%%   - Absorbs and supersedes: Surgical Derivation Note G1 ("The Sol
%%     Monodromy Forces the Hurwitz Enlargement of the Corner Order",
%%     corpus working note, June 2026). Hypothesis H1 [M] of that note
%%     is here replaced by the Transport Principle [D] and discharged
%%     in Section 8.
%%   - Incorporates: Deliverable G2 ("Canonical Construction of the
%%     Trace-phi Unit on M_12", June 2026).
%%   - Significance framing per PI adjudication (June 2026): structural
%%     rigidity (Tier 2), explicitly not empirical confirmation; the
%%     ring-theoretic/thermodynamic "order" pun is excluded from this
%%     paper (reserved for the popular treatment).
%%
%% House conventions: UK English; omega reserved corpus-wide for the
%% negentropic order parameter; Hurwitz unit written \varpi; golden
%% ratio written \varphi; rigour scale [R]/[D]/[M]/[G] on every step.
%%=====================================================================
\documentclass[11pt,a4paper]{article}

%% ---- Packages ----
\usepackage[utf8]{inputenc}
\usepackage[T1]{fontenc}
\usepackage[UKenglish]{babel}
\usepackage{amsmath,amssymb,amsthm,amsfonts}
\usepackage{mathtools}
\usepackage{geometry}
\usepackage{graphicx}
\usepackage{booktabs}
\usepackage{enumitem}
\usepackage{xcolor}
\usepackage{tcolorbox}
\usepackage{tabularx}
\usepackage{caption}
\usepackage{fancyhdr}
\usepackage{longtable}
\usepackage{array}
\usepackage{hyperref}

%% ---- Page Setup ----
\geometry{margin=2.5cm}
\hypersetup{
    colorlinks=true,
    linkcolor=blue!60!black,
    citecolor=green!50!black,
    urlcolor=blue!70!black
}

\pagestyle{fancy}
\fancyhf{}
\renewcommand{\headrulewidth}{0pt}
\fancyfoot[C]{\thepage}
\fancyfoot[R]{\small\textcopyright\ Dhiren J.\ Master 2026}
\fancypagestyle{plain}{\pagestyle{fancy}}

%% ---- tcolorbox styles ----
\tcbuselibrary{skins,breakable}
\newtcolorbox{statusbox}[1][]{
  enhanced, colback=green!3, colframe=green!40!black,
  fonttitle=\bfseries, breakable, title={#1}}
\newtcolorbox{definitionbox}[1][]{
  enhanced, colback=blue!2, colframe=blue!40!black,
  fonttitle=\bfseries, breakable, title={#1}}
\newtcolorbox{resultbox}[1][]{
  enhanced, colback=blue!4, colframe=blue!50!black,
  fonttitle=\bfseries, breakable, title={#1}}
\newtcolorbox{mechanismbox}[1][]{
  enhanced, colback=gray!4, colframe=gray!55!black,
  fonttitle=\bfseries, breakable, title={#1}}
\newtcolorbox{gapbox}[1][]{
  enhanced, colback=red!3, colframe=red!50!black,
  fonttitle=\bfseries, breakable, title={#1}}

%% ---- Theorem environments ----
\theoremstyle{plain}
\newtheorem{theorem}{Theorem}
\newtheorem{lemma}{Lemma}
\newtheorem{corollary}{Corollary}
\newtheorem{proposition}{Proposition}
\theoremstyle{definition}
\newtheorem{definition}{Definition}
\theoremstyle{remark}
\newtheorem{remark}{Remark}

%% Witness environment (W-1 ... W-4)
\newcounter{witness}
\newenvironment{witness}[1][]{\refstepcounter{witness}\par\medskip\noindent\textbf{W-\thewitness\ #1.}\itshape}{\par\medskip}

%% ---- Canonical MEF macro set (single source of truth) ----
%% Supplies \QCD (Quantum Coherence Dilaton, = \Phi_{QC}), \consciousness
%% (= \omega), \Mtwelve, \Keight, \C, \R, \Z, \Q (rationals, never the
%% dilaton), \Spinc, rigour marks, etc.  Do not redefine these locally.
\input{mef_macros}

%% ---- Paper-specific commands (not in mef_macros.tex) ----
\newcommand{\HH}{\mathbb{H}}
\newcommand{\Ico}{\mathbb{I}}
\newcommand{\Uset}{\mathcal{U}}
\newcommand{\Zphi}{\mathbb{Z}[\varphi]}
\newcommand{\FF}[1]{\mathbb{F}_{#1}}
\newcommand{\pfive}{\mathfrak{p}_5}
%% Framework order: calligraphic O ('order'), NOT \omega.  Macro named
%% \OrderMEF to avoid any visual or source-level confusion with the
%% consciousness coherence field \consciousness = \omega.
\newcommand{\OrderMEF}{\mathcal{O}_{\mathrm{MEF}}}
\newcommand{\conj}[1]{\operatorname{conj}_{#1}}
\newcommand{\rtag}[1]{\textnormal{\textbf{[#1]}}}
\DeclareMathOperator{\trd}{trd}
\DeclareMathOperator{\nrd}{nrd}
\DeclareMathOperator{\disc}{disc}
\DeclareMathOperator{\Sym}{Sym}
\DeclareMathOperator{\ord}{ord}
\DeclareMathOperator{\ImH}{Im}

%% ---- Title ----
\title{\textbf{The Icosian Completion of $\Keight$}\\[0.5em]
\large Sol Monodromy Forces the Hurwitz and Trace-$\varphi$ Enlargements:\\
$E_8$ as a Derived Arithmetic Lattice}
\author{D.~Master\thanks{\texttt{dhiren@masterequation.org}. \textcopyright\ Dhiren J.~Master 2026. UK English throughout.}}
\date{June 2026}

\begin{document}
\maketitle
\thispagestyle{fancy}

\begin{abstract}
\noindent
Within the Master Equation Framework on $\Mtwelve = (S^1_{\mathrm{Sol}} \times S^3) \times \Keight$, the pillowcase factor $T^2/\Z_2$ of $\Keight$ carries a quaternionic corner dictionary whose integral structure is, naively, the Lipschitz order. We prove that the framework's own Sol monodromy $M = \left(\begin{smallmatrix}2&1\\1&1\end{smallmatrix}\right)$, transported through this dictionary by a single Spin$^c$ Transport Principle, forces two successive enlargements of the corner order, one at each of the two arithmetically distinguished primes of the construction. The mod-$2$ shadow of $M$ (the orbifold prime) forces the Lipschitz $\to$ Hurwitz enlargement through the order-$3$ unit $\varpi = \tfrac12(-1+i+j+k)$; the mod-$\sqrt5$ shadow (the golden discriminant prime) forces the Hurwitz $\to$ icosian enlargement through a canonically constructed order-$10$ unit $q_5 = \tfrac12(\varphi + \varphi^{-1}i + j)$ of reduced trace $\varphi$, whose axis is the Sol unstable eigenvector, whose order is the order of $M$ in $GL(2,\FF{5})$, and whose half-integrality is the half-angle of the spin double cover. The framework order is therefore the icosian ring $\Ico$, the unique maximal order of Hamilton's quaternions over $\Q(\sqrt5)$, isometric to the $E_8$ root lattice under the Conway--Sloane rationalised norm. $E_8$ thus enters the framework as a \emph{derived arithmetic lattice} --- never as a gauge group --- with no continuous modulus surviving into the construction; in particular $\tau_2$ does not enter. The significance claimed is structural rigidity, not empirical confirmation.
\end{abstract}

\begin{statusbox}[Status, scope and provenance]
\textbf{Status:} standalone derivation paper within the Master Equation Framework (MEF) corpus; outside the numbered paper series by PI decision.

\smallskip
\textbf{Scope:} closes Gaps G1 and G2 of the icosian programme and discharges Hypothesis H1. Supersedes the corpus working note \emph{Surgical Derivation Note G1 --- The Sol Monodromy Forces the Hurwitz Enlargement of the Corner Order}, whose content is absorbed into \S\ref{sec:level2} with its Hypothesis H1 \rtag{M} replaced by the Transport Principle \rtag{D} of \S\ref{sec:tp} and discharged in \S\ref{sec:twoeval}.

\smallskip
\textbf{Rigour summary:} Lemmas 0--3 and Witnesses W-1--W-4 at \rtag{R}. Theorem~\ref{thm:hurwitz} (Hurwitz Forcing) and Theorem~\ref{thm:q5} (trace-$\varphi$ unit): implications at \rtag{R}, net classification \rtag{D} conditional on the Transport Principle, itself \rtag{D} (the single identified framework input; its level-2 instance is settled corpus data at \rtag{R}). Theorem~\ref{thm:icosian} at \rtag{R} given Theorems~\ref{thm:hurwitz}--\ref{thm:q5}. One convention at \rtag{M} (GR-1, consequence bounded). The Coldea remark (\S\ref{sec:coldea}) and the $E_4$ outlook (\S\ref{sec:outlook}) are observations and a programme target respectively, both \rtag{M}.

\smallskip
\textbf{Consumes (settled inputs, not re-derived):} Paper XV \S8 Theorem \texttt{thm:3cycle} (Sol Permutation Theorem) \rtag{R} \cite{PaperXV}; Paper XX $\Sigma$-2g/\,$\mathcal{M}$-7 (Sol monodromy data) \rtag{R} \cite{PaperXX}; Paper XXI \S4.1--4.4 (corner dictionary, $Q_8$ module action, Spin$^c$ reduction) \rtag{R/D} \cite{PaperXXI}; Paper III Addendum \S5 (corner assignment convention) \rtag{D} \cite{PaperIIIAdd}; Spin$^c$ Generation Halving Bridge ($J^2 = -1$) \cite{SpincBridge}.

\smallskip
\textbf{Feeds:} corpus-wide citations of $E_8$ as derived arithmetic structure; the $E_4$ theta-series programme target (\S\ref{sec:outlook}); the popular treatment \emph{The Shape of Becoming}.
\end{statusbox}

\paragraph{Symbol-collision note (mandatory, per corpus protocol).}
The letter $\omega$ is reserved corpus-wide for the negentropic order parameter (Paper XX, Axiom $\Sigma$-4) and does not appear in this paper. The Hurwitz unit conventionally denoted $\omega$ in the external literature (Conway--Sloane) is therefore written $\varpi$ throughout:
\begin{equation}\label{eq:varpi}
\varpi := \tfrac12(-1+i+j+k), \qquad
\varpi_0 := \tfrac12(1+i+j+k), \qquad
\varpi = \varpi_0^{\,2}, \quad \varpi^3 = 1.
\end{equation}
The golden ratio is written $\varphi$ (lower case) throughout. Per the canonical macro set (\texttt{mef\_macros.tex}), the Quantum Coherence Dilaton is denoted $\QCD$ and does not appear in this paper; $\Q$ denotes the field of rationals, never the dilaton. The framework order of Definition \ref{def:order} is written $\OrderMEF$ --- a calligraphic O, for ``order'' in the arithmetic sense --- and is unrelated to the consciousness coherence field $\consciousness$, which likewise does not appear here.

\tableofcontents

%%=====================================================================
\section{Introduction and motivation}\label{sec:intro}
%%=====================================================================

The corner dictionary of Paper XXI \cite{PaperXXI} assigns to the four corners of the pillowcase orbifold $T^2/\Z_2 \subset \Keight$ the four limiting sub-algebras $\{\R, \C_i, \C_j, \C_k\}$ of Hamilton's quaternion algebra, and the integral structure this generates is the Lipschitz order --- the naive integer quaternions, with norm-one unit group $Q_8$. Separately, the Sol monodromy $M = \left(\begin{smallmatrix}2&1\\1&1\end{smallmatrix}\right)$ acts on the pillowcase, and Paper XV \cite{PaperXV} proves \rtag{R} that this action fixes one corner and cyclically permutes the other three. This paper answers the question of whether these two settled facts are compatible \emph{at the level of the integral quaternion order} --- and shows that the incompatibility, pursued at the two arithmetically distinguished primes of the construction, forces the corner order all the way up to the icosian ring.

The result is a two-stage forcing chain:
\begin{equation}\label{eq:chain}
L \;\subset\; H \;\subset\; \Ico
\quad\Big\|\quad
Q_8 \;\subset\; 2T \;\subset\; 2I
\quad\Big\|\quad
\text{corners} \;\subset\; \text{corners} + \mathrm{Sol}\,(3) \;\subset\; \text{corners} + \mathrm{Sol}\,(3,5),
\end{equation}
in which each inclusion is driven by one reduction of the \emph{same} monodromy matrix: the Lipschitz $\to$ Hurwitz step by the mod-$2$ shadow of $M$ (the orbifold prime, at which the pillowcase corners live), and the Hurwitz $\to$ icosian step by the mod-$\pfive$ shadow (the ramified prime of the golden field, $\pfive = (\sqrt5)$, at which the Sol eigenvalues collide). Both steps are evaluations of a single Spin$^c$ Transport Principle (\S\ref{sec:tp}); the mod-$2$ evaluation produces the order-$3$ Hurwitz unit $\varpi$, and the mod-$\pfive$ evaluation canonically produces the order-$10$ trace-$\varphi$ unit
\begin{equation}\label{eq:q5}
q_5 \;=\; \tfrac12\big(\varphi + \varphi^{-1} i + j\big), \qquad
\nrd(q_5) = 1, \quad \trd(q_5) = \varphi, \quad q_5^{\,5} = -1 .
\end{equation}
Together they generate the binary icosahedral group $2I$, identifying the framework order as the icosian ring $\Ico$ --- the unique maximal order of Hamilton's quaternions over $\Q(\sqrt5)$ --- which is isometric to the $E_8$ root lattice under the Conway--Sloane rationalised norm \cite{ConwaySloane, Tits}.

\paragraph{One-line crux.}
\emph{A 3-cycle of the imaginary axes is invisible to every Lipschitz unit, whose conjugation action preserves every axis; by Skolem--Noether and the $SO(3)$ classification, any inner realisation is conjugation by a half-integer Hurwitz unit --- so the mod-2 Sol shadow puts $\varpi$ in the order. The mod-$\pfive$ shadow then repeats the mechanism one prime higher: the order-10 reduced monodromy, transported along the Sol unstable eigenvector through the spin double cover, is conjugation by the trace-$\varphi$ unit $q_5$ --- and $\langle 2T, q_5 \rangle = 2I$.}
The machinery is elementary throughout: two modular reductions, one $SO(3)$ classification, one half-angle, one norm argument, and the ADE classification of finite quaternion groups.

\paragraph{Significance claimed (stated up front).}
The significance of this result is calibrated explicitly in \S\ref{sec:significance} and is of the \emph{structural-rigidity} kind: independently motivated framework ingredients --- $M$ selected for Sol geometry and the Weinberg-angle route, the pillowcase for the generation structure, the quaternionic fibre by the Spin$^c$ dictionary --- converge without adjustment on a unique canonical arithmetic object, with each of the two distinguished primes supplying exactly one missing generator and nothing left over. This raises the cost of dismissing $\Keight$ as an arbitrary construction; it does not constitute physical confirmation, which in this corpus is carried solely by numerical predictions against measurement. Equally important is what the result is \emph{not}: $E_8$ enters as a derived arithmetic lattice, and the framework formally excludes $E_8$ as a fundamental gauge group, so the known obstructions to $E_8$ gauge unification do not apply (\S\ref{sec:notgauge}).

%%=====================================================================
\section{Setting}\label{sec:setting}
%%=====================================================================

All definitions are chained type-first; every symbol appearing in a numbered statement is defined here or in \eqref{eq:varpi}.

\subsection{Scalars}\label{sec:scalars}

$F := \Q(\sqrt5)$, a real quadratic number field; $\varphi := (1+\sqrt5)/2 \in F$, the golden ratio, satisfying $\varphi^2 = \varphi + 1$ (hence $\varphi^{-1} = \varphi - 1$ and $\varphi^{-2} = 2 - \varphi$); $\mathcal{O}_F = \Zphi$, its ring of integers; $\pfive := (\sqrt5) = (2\varphi - 1)$, the unique ramified prime of $F$, with residue field $\FF{5}$.

\setcounter{lemma}{-1}
\begin{lemma}[Sol eigenvalue ring is maximal]\label{lem:ring}
\rtag{R} $\Z[\varphi^{\pm2}] = \Zphi = \mathcal{O}_F$.
\end{lemma}
\begin{proof}
$\varphi = \varphi^2 - 1 \in \Z[\varphi^2]$, and $\varphi^{-2} = 2 - \varphi$.
\end{proof}

\noindent\emph{Gloss.} The ring generated by the Sol monodromy eigenvalues $\varphi^{\pm2}$ (Paper XX, $\mathcal{M}$-7 \cite{PaperXX}) is not a sub-ring of the golden ring: it \emph{is} the golden ring, the maximal order of $F$. Moreover $\varphi^2$ is the fundamental totally positive unit of $\mathcal{O}_F$ --- the Sol eigenvalues generate the totally positive unit group exactly.

\subsection{Algebras, orders, and norm-one unit groups}\label{sec:orders}

$B_{\Q} := \left(\frac{-1,\,-1}{\Q}\right)$, Hamilton's rational quaternion algebra, with standard basis $\{1,i,j,k\}$, reduced trace $\trd$, reduced norm $\nrd$, and imaginary part $\ImH\HH \cong \R^3$ at either real place; $B := B_{\Q} \otimes_{\Q} F$, totally definite, and unramified at every finite place of $F$ (the only candidate is $2$, which is inert in $F$ since $5 \equiv 5 \pmod 8$, and the unramified quadratic extension $\Q_2(\sqrt5)$ splits the local division algebra) \rtag{R}; consequently every maximal order of $B$ has reduced discriminant $(1)$, and the maximal order is unique up to conjugacy \cite{Vigneras, Tits}.

For an order $\mathcal{O}$ in $B$ (or $B_\Q$), write $\mathcal{O}^1 := \{u \in \mathcal{O} : \nrd(u) = 1\}$ for its \emph{norm-one unit group}, a finite group since $B$ is definite.\footnote{Earlier corpus working notes wrote $\mathcal{O}^\times$ for this finite group; the notation is refined here, since over $\Zphi$ the full unit group $\mathcal{O}^\times$ contains the infinite central group $\Zphi^\times$.} The orders in play:
\begin{itemize}[leftmargin=2em]
\item $L := \Z \oplus \Z i \oplus \Z j \oplus \Z k \subset B_\Q$ --- the \textbf{Lipschitz order}, the integral structure generated by the corner data of Paper XXI \S4.2. $L^1 = Q_8 = \{\pm1, \pm i, \pm j, \pm k\}$ \rtag{R} (the only integer solutions of $a^2+b^2+c^2+d^2 = 1$).
\item $H := \Z\text{-span}\langle 1, i, j, \varpi_0\rangle \subset B_\Q$ --- the \textbf{Hurwitz order} \cite{Hurwitz}; equivalently $H = L + \Z\varpi$. $H^1 = 2T \cong SL(2,3)$, the binary tetrahedral group, $|2T| = 24$, with $2T = Q_8 \rtimes \langle\varpi\rangle$ \rtag{R} (classical).
\item $H_\varphi := H \otimes_\Z \Zphi \subset B$ --- the golden-extended Hurwitz order, of reduced discriminant $2\Zphi \neq (1)$: strictly non-maximal \rtag{R}.
\item $\Ico := \Zphi\text{-span of } 2I \subset B$ --- the \textbf{icosian ring}, the unique maximal order of $B$ up to conjugacy, with $\Ico^1 = 2I$, the binary icosahedral group, $|2I| = 120$ \rtag{R, classical} \cite{ConwaySloane, Tits, Vigneras}.
\end{itemize}

\subsection{Geometry}\label{sec:geometry}

The arena is $\Mtwelve = (S^1_{\mathrm{Sol}} \times S^3) \times \Keight$ with $\Keight = [(\mathbb{CP}^2 \times S^2) \times_w (T^2/\Z_2)]^{\mathrm{Spin}^c}$, $\chi(\Keight) = 12$. The pillowcase is $T^2/\Z_2$, the quotient of $T^2 = \R^2/\Z^2$ by the involution $\sigma : x \mapsto -x$, with corners (the 2-torsion points)
\begin{equation}\label{eq:corners}
C_0 = (0,0), \qquad C_1 = (\tfrac12, 0), \qquad C_2 = (0, \tfrac12), \qquad C_3 = (\tfrac12, \tfrac12).
\end{equation}
The fibre over the pillowcase interior carries Hamilton's algebra $\HH$; the corner limits are the sub-algebras $\R, \C_i, \C_j, \C_k$, where $\C_i = \R \oplus \R i$ is the complex sub-algebra through the imaginary unit $i$, and likewise $\C_j, \C_k$. The \textbf{corner dictionary} is the map
\begin{equation}\label{eq:dictionary}
\mathcal{D} : \{C_0, C_1, C_2, C_3\} \longrightarrow \{\R, \C_j, \C_i, \C_k\}, \qquad
\mathcal{D}(C_0) = \R, \;\; \mathcal{D}(C_2) = \C_i, \;\; \mathcal{D}(C_1) = \C_j, \;\; \mathcal{D}(C_3) = \C_k,
\end{equation}
with the coordinate-to-algebra assignment fixed by the corpus corner architecture (Paper III Addendum \S5 \cite{PaperIIIAdd}; Paper XXI \S4.2). Half-period translations lift to right-multiplication by elements of $Q_8$ (Paper XXI \S4.3): the $x$ half-period (corner $C_1$) through $\pm j$, the $y$ half-period (corner $C_2$) through $\pm i$. The Spin$^c$ structure breaks $Sp(1) \to U(1)$ with determinant line $c_1(L) = 3$, selecting $\C_i$; the orbifold lift of $\sigma$ satisfies $J^2 = -1$ \cite{SpincBridge}; and the lowest Spin$^c$ Dirac eigenvalue on $T^2$ with $c_1(L) = 3$ is $\pi$, independent of the modulus $\tau_2$ (Paper XXI \S2.6).

\subsection{Dynamics}\label{sec:dynamics}

The Sol monodromy is $M = \left(\begin{smallmatrix}2&1\\1&1\end{smallmatrix}\right) \in SL(2,\Z)$, with characteristic polynomial $x^2 - 3x + 1$, eigenvalues $\varphi^{\pm2}$, eigenvectors $v_+ \propto (\varphi, 1)$ and $v_- \propto (1, -\varphi)$, and Weyl discriminant $5 = \disc\Q(\sqrt5)$ (Paper XX, $\Sigma$-2g and $\mathcal{M}$-7) \rtag{R}. It acts on $T^2$ linearly, commutes with $\sigma$, and hence descends to the pillowcase, preserving the corner set \eqref{eq:corners}. On the corner set, $M$ fixes $C_0$ and induces the 3-cycle $C_1 \to C_2 \to C_3$ (Sol Permutation Theorem, Paper XV \S8 \texttt{thm:3cycle}, settled \rtag{R}).

%%=====================================================================
\section{The Transport Principle}\label{sec:tp}
%%=====================================================================

The single framework input of this paper is stated before any result that uses it. It replaces, strengthens, and will retrospectively discharge (\S\ref{sec:twoeval}) the Hypothesis H1 of the superseded Note G1.

\begin{definition}[Frame map $\Psi$]\label{def:frame}
\rtag{D} Let $\Psi : \R^2 \to \ImH\HH$ be the $\R$-linear extension of the lattice-frame assignment fixed by the half-period lifts of \S\ref{sec:geometry}:
\begin{equation}\label{eq:frame}
\Psi(1,0) = j, \qquad \Psi(0,1) = i,
\end{equation}
defined up to the $Q_8$ sign projectivity of the lifts (lines $\R\,\Psi(v)$ are well defined). $\Psi$ uses only the affine lattice structure of $T^2$, not its conformal modulus: \emph{$\tau_2$ cannot enter through $\Psi$.}
\end{definition}

\begin{definitionbox}[Transport Principle (TP) --- the single framework input \;\rtag{D}]
The Sol monodromy acts on the torsion strata of the pillowcase fibre compatibly with the corner dictionary \eqref{eq:dictionary} and the frame \eqref{eq:frame}, as an $F$-algebra automorphism of the fibre algebra $\HH$, lifted through the Spin$^c$ double cover so that central parities are remembered. Explicitly, for each torsion level $N$ at which $M$ acts:

\smallskip
\textbf{(a) Equivariance.} The transport realises the action of $M$ on the level-$N$ torsion data of the pillowcase fibre, through $\mathcal{D}$ and $\Psi$, by an $F$-algebra automorphism $\alpha_M^{(N)}$ of $\HH$. By Skolem--Noether \rtag{R}, every such automorphism is inner: $\alpha_M^{(N)} = \conj{q_N}$ for a unit $q_N \in B^\times$, determined up to the centre $F^\times$.

\smallskip
\textbf{(b) Spinorial faithfulness.} The realisation factors through the spin double cover and carries the central sign of the reduced monodromy faithfully: if a power of $M$ reduces to $-I$ (respectively $+I$) on the level-$N$ torsion, the corresponding power of $q_N$ is the central element $-1$ (respectively $\pm1$, the sign being invisible exactly when $-I \equiv +I$ at that level).

\smallskip
\textbf{Settled instance.} The level-2 instance of (b) is existing corpus data \rtag{R}: $\sigma^2 = \mathrm{id}$ lifts to $J^2 = -1$ \cite{SpincBridge}, and the half-periods square to $-1$ inside $Q_8$ (Paper XXI \S4.3). The level-2 instance of (a) is the algebra-level reading of the Sol Permutation Theorem. The level-5 use in \S\ref{sec:level5} extends the same lattice frame $\R$-linearly; nothing else is assumed.
\end{definitionbox}

\begin{definition}[Framework order]\label{def:order}
$\OrderMEF \subset B$ is the $\Zphi$-order generated by the corner data (the Lipschitz generators of \S\ref{sec:orders}, extended by $\Zphi$ scalars) together with the transport realisations $\{q_N\}$ supplied by TP. The content of Theorems \ref{thm:hurwitz}, \ref{thm:q5} and \ref{thm:icosian} is the computation of $\OrderMEF$.
\end{definition}

\paragraph{Why this principle matters, and its status.}
TP is the exact bridge between two settled corpus layers. The \emph{sector-level} statement --- that the Sol action fixes the untwisted sector and 3-cycles the three twisted sectors --- is Paper XV Theorem \texttt{thm:3cycle}, at \rtag{R}. TP lifts this from the indexing set of sectors to the fibre algebra itself, with the double-cover clause (b) fixing how central signs travel. Its classification is \rtag{D}: it is the single identified framework input of the paper, with its level-2 evaluation already settled corpus data and its level-5 evaluation the new application; it is not derived here from deeper principles (Gap Register GR-3, where the bounded derivation target --- Spin$^c$ parallel transport of the fibre algebra around $S^1_{\mathrm{Sol}}$, patterned on the existing $V_{\mathrm{geom}}$ module-action lift of Paper XXI \S4.3 --- is located, not hidden). Everything new in this paper flows through TP; nothing else is assumed. The improvement over the superseded Note G1 is structural: where G1 assumed an integrality clause H1(b) \rtag{M}, here integrality is not assumed but \emph{checked} --- the realising units are constructed with explicit coefficients (Theorems \ref{thm:hurwitz} and \ref{thm:q5}), and the framework order is then computed from Definition \ref{def:order}.

%%=====================================================================
\section{Level-2 evaluation: the Hurwitz forcing}\label{sec:level2}
%%=====================================================================

The 2-torsion points of $T^2$ are precisely the corners \eqref{eq:corners}; the level-2 evaluation of TP is therefore the statement that the Sol corner 3-cycle acts inside the fibre algebra. This section shows that this requirement has exactly one integral cost: $\varpi$.

\begin{lemma}[Sol 3-cycle on corners --- consumed]\label{lem:3cycle}
\rtag{R} The descended action of $M$ on the corner set fixes $C_0$ and induces the 3-cycle $C_1 \to C_2 \to C_3 \to C_1$.
\end{lemma}

\noindent\emph{Source.} Paper XV \S8, Theorem \texttt{thm:3cycle} (Sol Permutation Theorem), proved there by direct evaluation on half-lattice coordinates; consumed here as a settled input.

\begin{remark}[One-line conceptual restatement]\label{rem:mod2}
Reduce mod 2: $\bar M = \left(\begin{smallmatrix}0&1\\1&1\end{smallmatrix}\right) \in GL(2,\FF2)$, identifying the corner set with $\FF2^2$ via \eqref{eq:corners}. The characteristic polynomial of $\bar M$ is $x^2 + x + 1$, irreducible over $\FF2$: $\bar M$ has order 3 and no non-zero fixed vector, hence acts in a single 3-cycle on the three non-zero points. \rtag{R} This places the geometric theorem of Paper XV inside $GL(2,\FF2) \cong S_3$ --- the arithmetic and geometric witnesses of the 3-cycle are the same object, and the mod-2 shadow of $M$ is the exact level-2 analogue of the mod-$\pfive$ shadow computed in \S\ref{sec:witnesses}.
\end{remark}

\begin{lemma}[Lipschitz units cannot see the 3-cycle]\label{lem:blind}
\rtag{R} The conjugation action of $L^1 = Q_8$ on the set of imaginary axes $\{\R i, \R j, \R k\}$ is trivial. In particular, no unit of the Lipschitz order realises any 3-cycle of the corner sub-algebras.
\end{lemma}
\begin{proof}
The centre $\{\pm1\}$ acts trivially. Conjugation by $i$ sends $(i,j,k) \mapsto (i,-j,-k)$: it preserves each axis as a line, merely reversing orientation on two of them. By the symmetry of Hamilton's relations the same holds for $\pm j, \pm k$. The induced map $Q_8 \to \Sym\{\R i, \R j, \R k\}$ is therefore the trivial homomorphism.
\end{proof}

\noindent\emph{Gloss.} This is the precise content of Paper XXI \S4.3's group $V_{\mathrm{Aut}}$ --- ``rotations by $\pi$ about the $i$- and $j$-axes'' --- read as a permutation statement: $V_{\mathrm{Aut}}$ flips signs along axes; it never exchanges axes. The Sol 3-cycle of Lemma \ref{lem:3cycle} is therefore not merely \emph{unrealised} by the corner unit group; it is \emph{unrealisable} by it. This is the obstruction whose unique resolution is Lemma \ref{lem:rigidity}.

\begin{lemma}[Rigidity: every inner realisation of a corner 3-cycle is Hurwitz, and half-integral]\label{lem:rigidity}
\rtag{R} Let $q \in B^\times$ (or $B_\Q^\times$) be such that $\conj{q} : x \mapsto qxq^{-1}$ permutes the three sub-algebras $\{\C_i, \C_j, \C_k\}$ in a 3-cycle. Then
\begin{equation}\label{eq:Uset}
q \in F^\times \cdot \Uset, \qquad
\Uset := \big\{ \tfrac12(\varepsilon_0 + \varepsilon_1 i + \varepsilon_2 j + \varepsilon_3 k) \;:\; \varepsilon_n \in \{\pm1\} \big\},
\end{equation}
the set of sixteen half-integer Hurwitz units. Every $u \in \Uset$ satisfies $\nrd(u) = 1$, $\trd(u) = \pm1$, and $L + \Z u = H$.
\end{lemma}
\begin{proof}
Every $F$-algebra automorphism of the quaternion algebra is inner (Skolem--Noether), and $\conj{q}$ depends only on the class of $q$ in $B^\times/F^\times$; the induced action on $\ImH\HH \cong \R^3$ (at either real place) is a rotation in $SO(3)$. A rotation permuting the three coordinate lines lies in the octahedral group; the elements of the octahedral group inducing a 3-cycle on the lines are exactly the eight rotations by $\pm2\pi/3$ about the body diagonals $(\pm1,\pm1,\pm1)/\sqrt3$. The unit-quaternion lifts of the rotation by $2\pi/3$ about $(\varepsilon_1,\varepsilon_2,\varepsilon_3)/\sqrt3$ are
\begin{equation}\label{eq:lifts}
\pm\Big( \cos\tfrac{\pi}{3} + \sin\tfrac{\pi}{3}\cdot\tfrac{\varepsilon_1 i + \varepsilon_2 j + \varepsilon_3 k}{\sqrt3} \Big)
= \pm\tfrac12\big(1 + \varepsilon_1 i + \varepsilon_2 j + \varepsilon_3 k\big),
\end{equation}
and the inverse rotations supply real part $-\tfrac12$; together these are exactly $\Uset$. The norm and trace values are immediate; and for any $u \in \Uset$, the difference $u - \varpi_0$ lies in $L$ (each coordinate discrepancy is an integer), so $L + \Z u = L + \Z\varpi_0 = H$.
\end{proof}

\noindent\emph{Gloss of the bare identities \eqref{eq:Uset}--\eqref{eq:lifts}.} The classification is sign-convention-free: whatever orientation the 3-cycle carries, and whatever signs the dictionary attaches to the imaginary units at each corner, \emph{every} inner realisation is a central scalar times a \textbf{half-integer} unit, and \emph{any single one} of the sixteen already generates the full Hurwitz order over the Lipschitz order. There is no realisation that stays integral in the naive sense. The half-integrality is forced, not chosen --- and its geometric source is identified in Theorem \ref{thm:q5}, Step 4: it is the half-angle of the spin double cover.

\begin{theorem}[Hurwitz Forcing]\label{thm:hurwitz}
Under TP, the level-2 evaluation $q_2$ satisfies:
\begin{enumerate}[label=(\roman*), leftmargin=2.2em]
\item $q_2 = \lambda u_0$ with $\lambda \in \Zphi^\times$ central and $u_0 \in \Uset$; hence $u_0 = \lambda^{-1} q_2 \in \OrderMEF$ and $\OrderMEF \supseteq H_\varphi$.
\item The norm-one unit group enlarges minimally and exactly: $Q_8 \subsetneq \langle Q_8, u_0\rangle = 2T$, the binary tetrahedral group, with the new order-3 generator supplied by the Sol realisation.
\end{enumerate}
\emph{Rigour:} the implication TP $\Rightarrow$ (i)--(ii) is \rtag{R}; the net classification is \rtag{D} conditional on TP \rtag{D}.
\end{theorem}
\begin{proof}
By TP(a) and Lemma \ref{lem:3cycle}, $\conj{q_2}$ induces a 3-cycle of $\{\C_i, \C_j, \C_k\}$. By Lemma \ref{lem:blind}, $q_2 \notin F^\times \cdot L^1$. By Lemma \ref{lem:rigidity}, $q_2 = \lambda u_0$ with $\lambda$ central and $u_0 \in \Uset$. Since $B$ is definite and $q_2$ is integral over $\Zphi$ (its explicit class is exhibited in \S\ref{sec:twoeval}), $\nrd(q_2) = \lambda^2 \in \Zphi^\times$ forces $\lambda \in \Zphi^\times$, and $u_0 = \lambda^{-1} q_2 \in \OrderMEF$ (orders are closed under the standard involution since $\trd(\OrderMEF) \subseteq \mathcal{O}_F$, so inverses of units stay in the order). Lemma \ref{lem:rigidity} then gives $L + \Z u_0 = H$, hence $H_\varphi \subseteq \OrderMEF$. For (ii): $\langle Q_8, u_0\rangle \subseteq H^1 = 2T$; the image of $u_0$ in $2T/\{\pm1\} \cong A_4$ has order 3 (Lemma \ref{lem:rigidity}: $\trd(u_0) = \pm1$, so $u_0$ has order 3 or 6, in either case order 3 modulo centre) and $u_0$ normalises $Q_8$, so $|\langle Q_8, u_0\rangle| \geq 3|Q_8| = 24 = |2T|$; equality follows.
\end{proof}

\noindent\emph{Prose interpretation.} The theorem converts a dynamical statement (the Sol clock permutes three corners) into an arithmetic one (the integral fibre algebra contains half-integer quaternions). The corner order cannot remain Lipschitz in a Sol-equivariant universe.

\begin{corollary}[Strict non-maximality; the residual gap is exactly trace-$\varphi$]\label{cor:nonmax}
\rtag{R} $H_\varphi$ has reduced discriminant $2\Zphi \neq (1)$, while every maximal order of $B$ has reduced discriminant $(1)$ (\S\ref{sec:orders}). Hence $H_\varphi$ is strictly non-maximal, and the residual enlargement $H_\varphi \subsetneq \Ico$ to the unique maximal order is exactly the adjunction of a finite-order unit of reduced trace $\varphi$ (order 10) or $\varphi - 1$ (order 5) --- impossible over $\Z$-traces, by the crystallographic restriction: finite-order units with trace in $\Z$ have orders in $\{1,2,3,4,6\}$, while orders 5 and 10 require trace in $\{\pm\varphi, \pm(\varphi-1)\}$, available only in $\Zphi$ --- the ring Lemma \ref{lem:ring} shows the Sol eigenvalues already generate.
\end{corollary}

\noindent\emph{Gloss.} The division of labour is exact: Theorem \ref{thm:hurwitz} exhausts what is achievable over $\Z$-traces; the missing generator must carry a golden trace, and the only golden structure in the framework is the Sol eigendata. This is the hand-off to the level-5 evaluation.

%%=====================================================================
\section{Arithmetic witnesses at the ramified prime}\label{sec:witnesses}
%%=====================================================================

Before constructing the trace-$\varphi$ unit, we verify four elementary statements about the mod-$\pfive$ shadow of $M$. Each is stated with a one-line proof; each is at \rtag{R}.

\setcounter{witness}{0}
\begin{witness}[{\rtag{R}}]
$M$ mod 5 has the single eigenvalue $4 = -1$ with a non-trivial Jordan block; hence $M$ has order 10 in $GL(2,\FF5)$, and $M^5 \equiv -I \pmod 5$.
\end{witness}
\noindent\emph{Proof.} $\chi_M(x) = x^2 - 3x + 1 \equiv (x-4)^2 \pmod 5$ (discriminant $9 - 4 = 5 \equiv 0$); $M - 4I \equiv \left(\begin{smallmatrix}3&1\\1&2\end{smallmatrix}\right) \pmod 5$ has determinant $5 \equiv 0$ and is non-zero, hence rank 1, so the Jordan block is non-trivial; in Jordan form $M \equiv \left(\begin{smallmatrix}-1&1\\0&-1\end{smallmatrix}\right)$, whence the unipotent part has order 5 and $M^5 \equiv -I$, giving $\ord(M) = 10$. (Directly: $M^5 = \left(\begin{smallmatrix}89&55\\55&34\end{smallmatrix}\right) \equiv \left(\begin{smallmatrix}4&0\\0&4\end{smallmatrix}\right) \pmod 5$, the Fibonacci entries $F_{11}, F_{10}, F_9$.) $\square$

\begin{witness}[{\rtag{R}}]
At the ramified prime $\pfive$ of $\Zphi$, the two Sol eigenvalues $\varphi^{\pm2}$ collide (both $\equiv 4 \bmod \pfive$): the stable and unstable directions degenerate, and the residual unipotent supplies the 5-fold structure.
\end{witness}
\noindent\emph{Proof.} $\sqrt5 = 2\varphi - 1 \equiv 0$, so $\varphi \equiv 3 \pmod{\pfive}$ (residue field $\FF5$, $2 \cdot 3 = 6 \equiv 1$); then $\varphi^2 \equiv 9 \equiv 4$ and $\varphi^{-2} \equiv 4^{-1} \equiv 4$. $\square$

\begin{witness}[{\rtag{R}}]
The pillowcase involution coincides with the fifth power of the Sol monodromy on 5-torsion: $\sigma = M^5$ on $T^2[5]$.
\end{witness}
\noindent\emph{Proof.} $\sigma = -\mathrm{id}$ acts on $T^2[5]$ as $-I$; by W-1, $M^5 \equiv -I$ on $T^2[5]$. $\square$

\begin{witness}[{\rtag{R}}]
The 24 non-zero 5-torsion points of $T^2$ descend to 12 pillowcase points ($\sigma$ pairs them freely), and $M$ acts on these 12 points with two fixed points and two 5-cycles; the fixed points lie on the mod-5 eigenline $\ker(M - 4I) = \mathrm{span}\{(1,2)\}$.
\end{witness}
\noindent\emph{Proof.} $\sigma$ is free on $T^2[5]\smallsetminus\{0\}$ ($x = -x \Rightarrow 2x = 0 \Rightarrow x = 0$, as $\gcd(2,5) = 1$), so $24/2 = 12$ pillowcase points. A pillowcase point is $M$-fixed iff $Mx = \pm x$; since $\chi_M(1) = -1 \not\equiv 0 \pmod 5$, only the eigenvalue $-1 \equiv 4$ contributes, with $\ker(M - 4I) = \mathrm{span}\{(1,2)\}$ (from $3a + b \equiv 0$), giving the two fixed classes $\{\pm(1,2)\}$ and $\{\pm(2,4)\}$. On the quotient $M^5 = \sigma = \mathrm{id}$ (W-3), so the remaining 10 points fall into two 5-cycles, explicitly
$\{\pm(1,0)\} \to \{\pm(2,1)\} \to \{\pm(0,3)\} \to \{\pm(3,3)\} \to \{\pm(4,1)\} \to$ and
$\{\pm(0,1)\} \to \{\pm(1,1)\} \to \{\pm(3,2)\} \to \{\pm(3,0)\} \to \{\pm(1,3)\} \to$. $\square$

\begin{remark}[Numerology guard]\label{rem:guard}
The count $12 = (5^2 - 1)/2$ of pillowcase 5-torsion points equals $\chi(\Keight)$. This is reported as a coincidence; no step in this paper uses it, and it is logged as Gap Register item GR-4 \rtag{G, observation only}.
\end{remark}

%%=====================================================================
\section{Level-5 evaluation: the canonical trace-\texorpdfstring{$\varphi$}{phi} unit}\label{sec:level5}
%%=====================================================================

The level-5 evaluation of TP must realise the order-10 action of W-1/W-4 inside the fibre algebra. This section constructs the realising unit canonically --- the deliverable is the mechanism, not the specimen --- and verifies that no continuous modulus survives into it.

\begin{theorem}[Canonical trace-$\varphi$ unit]\label{thm:q5}
Let $\bar\ell = \ker(M - 4I) \subset \FF5^2$ (W-4), and impose on the level-5 evaluation $q$ of the transport:
\begin{enumerate}[label=\textup{(T\arabic*)}, leftmargin=2.6em]
\item \textbf{axis condition:} conjugation by $q$ fixes the $\Psi$-image of the $M$-invariant $\Zphi$-line lifting $\bar\ell$, selected by Sol instability;
\item \textbf{spinorial faithfulness:} $\ord(q) = \ord(M \bmod \pfive) = 10$ with $q^5 = -1$, mirroring $M^5 \equiv -I$ (W-1), per TP(b);
\item \textbf{step normalisation:} conjugation by $q$ advances each 5-cycle of W-4 by one step, realised as the minimal rotation $2\pi/5$ \emph{(convention --- Gap Register GR-1)}.
\end{enumerate}
Then $q$ is, up to $2T$-conjugacy and inversion (equivalently, up to the $Q_8$ sign ambiguities of $\Psi$ and the orientation of $v_+$),
\begin{equation}\label{eq:q5main}
q_5 = \tfrac12\big(\varphi + \varphi^{-1} i + j\big), \qquad
\nrd(q_5) = 1, \quad \trd(q_5) = \varphi, \quad q_5^{\,5} = -1,
\end{equation}
and no continuous modulus --- in particular $\tau_2$ --- enters the construction.
\emph{Rigour:} Steps 1, 2 \rtag{D} given TP (with \rtag{R} sub-steps as marked); Step 3 \rtag{M} for the convention, \rtag{R} thereafter; Steps 4, 5 \rtag{R}.
\end{theorem}

\begin{proof}
\textbf{Step 1 --- Axis \rtag{D, rigorous given TP}.}
By W-4 the $M$-fixed pillowcase 5-torsion locus spans $\bar\ell = \mathrm{span}_{\FF5}\{(1,2)\}$. Equivariance (TP(a)) forces the realising rotation to fix the $\Psi$-image of the fixed locus, so the axis is $\Psi$ of a canonical lift of $\bar\ell$. Over $F$, $M$ has distinct eigenvalues $\varphi^{\pm2}$ \rtag{R}, so its invariant lines are exactly the eigenlines $\R v_+, \R v_-$. Both reduce to $\bar\ell$: $v_+ = (\varphi, 1) \equiv (3,1) \propto (1,2)$ and $v_- = (1, -\varphi) \equiv (1,2) \bmod \pfive$ --- this is precisely the W-2 collision, so the mod-$\pfive$ axis condition determines the axis only up to the stable/unstable dichotomy. The Sol dynamics canonically distinguishes the \textbf{unstable} line (eigenvalue $\varphi^2 > 1$), given the orientation of $S^1_{\mathrm{Sol}}$ (GR-2). Hence
\begin{equation}\label{eq:axis}
\text{axis } a = \Psi(v_+) = i + \varphi j
\end{equation}
--- the golden weights enter here, sourced entirely in the Sol eigenbasis. \emph{(Deliverable (ii) of the G2 brief discharged.)}

\medskip
\textbf{Step 2 --- Order and central parity \rtag{D; rationality sub-step R}.}
Conjugation by $q$ acts on the 12-point locus with order 5 (W-4: on the pillowcase quotient, $M^5 = \sigma = \mathrm{id}$). Therefore $q^5$ acts trivially by conjugation on the locus, so $q^5$ is central in the algebra generated --- unless $q^5$ were a non-central unit with trivial 12-point action; but any such unit shares $q$'s axis, and the only candidates $\pm J = \pm k$ (the frame lift of $\sigma$) have axis $k \perp a$, while an order-20 unit would require $\trd = 2\cos(\pi/10) = \sqrt{(5+\sqrt5)/2} \notin F$ (its square $\varphi\sqrt5$ has norm $5$, not a square in $\Q$): both excluded \rtag{R}. Hence $q^5 \in \{\pm1\}$. Spinorial faithfulness (TP(b): the double cover remembers parities, exactly as $J^2 = -1$ remembers $\sigma^2 = \mathrm{id}$) sends the non-trivial central shadow $M^5 \equiv -I$ (W-1) to the non-trivial central element: $q^5 = -1$, so $\ord(q) = 10$. The eigenvalues of $q$ are primitive 10th roots $e^{\pm i\alpha}$, $\alpha \in \{\pi/5, 3\pi/5\}$, so the rotation angle is $2\pi/5$ or $4\pi/5$ and $\trd(q) = 2\cos\alpha \in \{\varphi, 1-\varphi\}$ \rtag{R}. \emph{(Deliverable (iii) discharged: the order 10 originates in W-1, the ramified-prime shadow --- in exact parallel with the mod-2 origin of $\varpi$'s order 3, Remark \ref{rem:mod2}.)}

\medskip
\textbf{Step 3 --- Angle \rtag{M for the convention, R thereafter}.}
The one-step normalisation (T3) identifies $M$'s unipotent step (W-1) with the minimal rotation $2\pi/5$, giving $\cos\alpha = \cos(\pi/5) = \varphi/2$, hence $\trd(q) = \varphi$. The alternative identification ($4\pi/5$) yields $q_5^{\,3}$ ($\trd = 1 - \varphi$), with $\langle q_5^{\,3}\rangle = \langle q_5\rangle$, so everything from Step 5 onward is invariant (GR-1).

\medskip
\textbf{Step 4 --- Explicit form; the source of the $\tfrac12$ \rtag{R}.}
The spin lift of the rotation by $2\pi/5$ about $\hat a = a/|a|$ is $q = \cos(\pi/5) + \sin(\pi/5)\,\hat a$. The half-angle of the double cover is the \emph{sole source of the factor $\tfrac12$}: $\cos(\pi/5) = \varphi/2$, and the two candidate sources named in the G2 brief (the half-period corner lattice; $J^2 = -1$) are the \emph{same} phenomenon --- both are instantiations of the binary lift. \emph{(Deliverable (i) discharged.)} For the imaginary part: $|a|^2 = 1 + \varphi^2 = 2 + \varphi$, $\sin(\pi/5) = \tfrac12\sqrt{3 - \varphi}$, and the golden identity $\varphi^2 = \varphi + 1$ gives $(3-\varphi) = (2-\varphi)(2+\varphi) = \varphi^{-2}(2+\varphi)$, whence $\sin(\pi/5)/|a| = \tfrac12\varphi^{-1}$. Therefore
\begin{equation}
q_5 = \tfrac{\varphi}{2} + \tfrac12\varphi^{-1}\,(i + \varphi j) = \tfrac12\big(\varphi + \varphi^{-1} i + j\big).
\end{equation}
Integrality in $\tfrac12\Zphi$ is not generic: it closes on $\varphi^2 + \varphi^{-2} = \operatorname{tr}(M) = 3$, i.e.\ on the Sol monodromy trace itself (equivalently $(\varphi + \varphi^{-1})^2 = 5 = \disc\Q(\sqrt5)$). No continuous modulus appears at any stage: the axis is projective-arithmetic, the angle discrete, the half topological. \textbf{Canonicality criterion passed; $\tau_2$ does not enter.}

\medskip
\textbf{Step 5 --- Well-definedness class \rtag{R}.}
The residual ambiguities are: (a) the $Q_8$ signs in $\Psi$, which flip axis components, moving the axis within the line set $\{(1, \pm\varphi, 0)\} \cup \text{cyclic}$ --- realised by conjugation inside $Q_8 \subset 2T$ (e.g.\ $\conj{i}$ maps the axis $(1,\varphi,0)$ to $(1,-\varphi,0)$); (b) the sign of $v_+$, which reverses the axis: $k\, q_5\, k^{-1} = \tfrac12(\varphi - \varphi^{-1}i - j) = q_5^{-1}$, so this is the inversion ambiguity; (c) the central twist $\pm q_5$, fixed by Step 2's parity. Hence $q_5$ is well defined up to $2T$-conjugacy and inversion, as claimed.
\end{proof}

\subsection*{Verification}
\begin{enumerate}[label=\textup{(V\arabic*)}, leftmargin=2.6em]
\item \textbf{Norm \rtag{R}:} $\nrd(q_5) = \tfrac14(\varphi^2 + \varphi^{-2} + 1) = \tfrac14(3 + 1) = 1$, using $\varphi^2 + \varphi^{-2} = (2\varphi - 1)^2 - 2 = 3 = \operatorname{tr} M$.
\item \textbf{Trace \rtag{R}:} $\trd(q_5) = \varphi$, as required by Corollary \ref{cor:nonmax}; the companion value $\varphi - 1 = \trd(q_5^{\,2})$ (order 5) is the alternative adjunction (GR-5).
\item \textbf{Order \rtag{R}:} the minimal polynomial is $q_5^{\,2} = \varphi q_5 - 1$, whence $q_5^{\,3} = \varphi(q_5 - 1)$ and $q_5^{\,5} = (\varphi q_5 - 1)\cdot\varphi(q_5 - 1) = \varphi(1 - \varphi) = -1$; $\ord(q_5) = 10 = \ord(M \bmod \pfive)$.
\end{enumerate}

%%=====================================================================
\section{The icosian completion}\label{sec:icosian}
%%=====================================================================

\begin{theorem}[Icosian completion; arithmetic $E_8$]\label{thm:icosian}
\rtag{R given Theorems \ref{thm:hurwitz}--\ref{thm:q5}} With $u_0$ ($=\varpi$ up to class, \S\ref{sec:twoeval}) and $q_5$ as constructed:
\begin{enumerate}[label=(\roman*), leftmargin=2.2em]
\item $q_5 \in \Ico$, and $\langle 2T, q_5\rangle = 2I$;
\item $\OrderMEF = \Ico$, the icosian ring, the unique maximal order of $B$ up to conjugacy;
\item under the Conway--Sloane rationalised norm, $\Ico$ is isometric to the $E_8$ root lattice \emph{(classical input, cite only \cite{ConwaySloane, Tits, MoodyPatera})} --- a derived arithmetic lattice, not a gauge group.
\end{enumerate}
\end{theorem}
\begin{proof}
(i) \emph{Membership:} the coefficient vector of $q_5$ in the basis $(1,i,j,k)$ is $\tfrac12(\varphi, \varphi^{-1}, 1, 0)$, the even coordinate permutation $(1\,4)(2\,3)$ of the standard icosian unit pattern $\tfrac12(0, 1, \varphi^{-1}, \varphi)$ \cite[ch.~8 \S2.1]{ConwaySloane}; hence $q_5 \in \Ico^1 = 2I$. (The stable-branch unit of Step 1, with coefficient pattern $\tfrac12(\varphi, \pm1, \mp\varphi^{-1}, 0)$, is the \emph{odd} permutation $(1\,4)$ of the same pattern, hence lies not in this copy of $\Ico$ but in a conjugate maximal order --- confirming GR-2: orientation reversal of $S^1_{\mathrm{Sol}}$ lands in the conjugate order, and the conclusion ``$\OrderMEF = \Ico$ up to conjugacy'' is orientation-independent.) \emph{Generation:} $2T \subset 2I$ and $\ord(q_5) = 10$, so $|\langle 2T, q_5\rangle|$ is divisible by $\operatorname{lcm}(24, 10) = 120 = |2I|$; hence $\langle 2T, q_5\rangle = 2I$. Structural confirmation via the ADE route: a finite subgroup of unit quaternions is cyclic, binary dihedral, $2T$, $2O$ or $2I$; a finite subgroup properly containing $2T$ together with an order-5 element must be $2I$ ($2O$ contains no order-5 element); over $F$ this is further strengthened arithmetically --- $2O$ is impossible outright, since its order-8 elements require $\trd = \pm\sqrt2 \notin \Q(\sqrt5)$, with $2$ inert in $\Zphi$ \rtag{R}.

(ii) Every generator of $\OrderMEF$ (the Lipschitz generators, $u_0$, $q_5$) lies in $\Ico$, and $\Ico$ is a ring, so $\OrderMEF \subseteq \Ico$. Conversely $\OrderMEF$ contains the group $2I$ by (i), hence its $\Zphi$-span, which is $\Ico$ \cite{ConwaySloane}. So $\OrderMEF = \Ico$; maximality and uniqueness up to conjugacy are classical \cite{Vigneras, Tits}. Consistency with Corollary \ref{cor:nonmax}: $\trd(q_5) = \varphi \notin \Z$ --- the enlargement is exactly the golden-trace adjunction the crystallographic restriction demanded, supplied by the Sol eigenbasis.

(iii) Cite-only, as stated.
\end{proof}

\subsection{The $(2,3,5)$ decomposition across the corpus}\label{sec:235}

The icosahedral triangle data $(2,3,5)$ --- equivalently, the element orders of $2I$ modulo centre --- decomposes across the corpus with no overlaps and no gaps:

\begin{center}
\begin{tabular}{@{}llll@{}}
\toprule
\textbf{Order(s)} & \textbf{Geometric source} & \textbf{Algebraic output} & \textbf{Site} \\
\midrule
2, 4 & pillowcase corners (2-torsion) & $Q_8 \subset L^1$ & Paper XXI \S4.2--4.3 \\
3 & Sol shadow at the orbifold prime $2$ & $\varpi$; $H^1 = 2T$ & this paper, \S\ref{sec:level2}, \S\ref{sec:twoeval} \\
5, 10 & Sol shadow at the golden prime $\pfive$ & $q_5$; $\Ico^1 = 2I$ & this paper, \S\ref{sec:level5} \\
\bottomrule
\end{tabular}
\end{center}

\begin{mechanismbox}[Construction summary: input data $\to$ mechanism $\to$ object identified]
\begin{center}
\small
\begin{tabularx}{\textwidth}{@{}>{\raggedright\arraybackslash}X >{\raggedright\arraybackslash}X >{\raggedright\arraybackslash}X@{}}
\toprule
\textbf{Input data} & \textbf{Mechanism} & \textbf{Output} \\
\midrule
Corner dictionary $\mathcal{D}$, half-period lifts (Paper XXI); Sol 3-cycle (Paper XV) & level-2 transport; Skolem--Noether; $SO(3)$ classification (Lemma \ref{lem:rigidity}) & $\varpi$; $L \to H$; $Q_8 \to 2T$ \\
\addlinespace
W-4 fixed locus $\bar\ell = \mathrm{span}\{(1,2)\}$; W-2 collision; Sol eigenbasis $v_+ \propto (\varphi,1)$; frame $\Psi$ & axis $= \Psi(\text{unstable lift of } \bar\ell)$ & $a = i + \varphi j$ (golden weights) \\
\addlinespace
W-1: $\ord(M \bmod \pfive) = 10$, $M^5 \equiv -I$; Spin$^c$ parity ($J^2 = -1$) & faithful binary lift; $q^5 = -1$ & $\ord(q_5) = 10$, $\trd(q_5) = \varphi$ \\
\addlinespace
double cover half-angle & $\cos(\pi/5) = \varphi/2$ & the factor $\tfrac12$ \\
\addlinespace
$\varphi^2 = \varphi + 1$; $\operatorname{tr} M = 3$ & integrality closure & $q_5 = \tfrac12(\varphi + \varphi^{-1}i + j) \in \Ico$ \\
\addlinespace
$\operatorname{lcm}(24,10) = 120$; ADE; $2$ inert excludes $2O$ & generation & $\langle 2T, q_5\rangle = 2I$; $\OrderMEF = \Ico \cong E_8$ (arithmetic, cite-only) \\
\bottomrule
\end{tabularx}
\end{center}
\end{mechanismbox}

%%=====================================================================
\section{One transport, two evaluations: H1 discharged}\label{sec:twoeval}
%%=====================================================================

The superseded Note G1 rested on a Hypothesis H1 \rtag{M}: that the Sol corner 3-cycle is realised by conjugation by a unit of the framework order. Here that realisation is \emph{constructed}, as the level-2 evaluation of the same transport whose level-5 evaluation produced $q_5$.

Evaluate the transport at the 2-torsion stratum. By Remark \ref{rem:mod2}, $M \bmod 2$ has order 3 and induces the corner 3-cycle $C_1 \to C_2 \to C_3$, i.e.\ (through $\mathcal{D}$) the sub-algebra cycle
\begin{equation}\label{eq:cycle}
\C_j \to \C_i, \qquad \C_i \to \C_k, \qquad \C_k \to \C_j,
\end{equation}
the unsigned axis cycle $(i\;k\;j)$. By TP(a) and Skolem--Noether, the realisation is conjugation by a unit inducing this line 3-cycle: an order-3 rotation about a body diagonal (Lemma \ref{lem:rigidity}). Exactly four rotations induce the given cyclic order: the order-3 rotations of the tetrahedral group $A_4$ fall into two conjugacy classes of four, one per cyclic order of the axes (the image of $A_4$ in $\Sym\{\text{axes}\} \cong S_3$ is $A_3$, with kernel the Klein group $V_4$). Their spin lifts therefore form a single $2T$-conjugacy class up to central sign --- \textbf{the class of $\varpi$}. Direct check \rtag{R}: with $\varpi$ as in \eqref{eq:varpi}, $\conj{\varpi_0}$ realises $(i\;j\;k)$ ($\varpi_0\, i\, \varpi_0^{-1} = j$), so $\conj{\varpi} = \conj{\varpi_0^2}$ realises $(i\;k\;j)$ --- precisely \eqref{eq:cycle}, matching the order-3 structure of Paper XV \texttt{thm:order3} on the nose \rtag{D conditional on the assignment convention of Paper III Addendum \S5; Theorem \ref{thm:hurwitz} does not depend on this match}.

The central parities are two evaluations of one rule, TP(b): at $p = 2$ the central sign is \emph{invisible} ($-I \equiv I \bmod 2$, and $M^3 \equiv I$), which is why $\varpi$ carries an irreducible $\pm$ ambiguity ($\varpi^3 = 1$, $(-\varpi)^3 = -1$, both compatible); at $\pfive$ the sign is \emph{visible} ($M^5 \equiv -I$, W-1) and forces $q_5^{\,5} = -1$ exactly.

\begin{proposition}[H1 discharged]\label{prop:h1}
\rtag{D, conditional only on TP} The corner 3-cycle is realised by conjugation by a constructed --- not hypothesised --- unit, $\pm\varpi \in \Uset \subset H \subset \Ico = \OrderMEF$, a unit of the framework order. The Hypothesis H1 of the superseded Note G1 thereby upgrades from \rtag{M} to \rtag{D} conditional on TP, and the net classification of Theorem \ref{thm:hurwitz} improves accordingly from ``\rtag{D} conditional on H1 \rtag{M}'' to ``\rtag{D} conditional on TP \rtag{D}''.
\end{proposition}

%%=====================================================================
\section{Significance: structural rigidity, not confirmation}\label{sec:significance}
%%=====================================================================

\subsection{Calibration of the claim}\label{sec:calibration}

Three kinds of evidence must be kept distinct. \emph{Empirical} evidence is numerical agreement with measurement; within this corpus it is carried by the derived observables ($\alpha^{-1}$, $\sin^2\theta_W$, $\Lambda_{\mathrm{eff}}$, $\eta$, $\Omega_{\mathrm{DM}}/\Omega_b$, \dots) and by nothing else. \emph{Structural-rigidity} evidence is the convergence of independently motivated ingredients, without adjustment, on a unique canonical mathematical object. \emph{Methodological} commitments --- here, the geometrisation wager in the Einstein--Kaluza--Klein lineage --- are vindicated only through outputs of the first kind. The present result belongs squarely to the second class, and the claim is calibrated accordingly.

The deflationary reading is stated openly, because it is true: the quaternion algebra $B$ over $\Q(\sqrt5)$ has exactly one maximal order up to conjugacy --- the icosian ring --- and the icosians are $E_8$. Once the framework contained a quaternionic fibre and golden-ratio scalars, $\Ico$ was the only place upward to go; in that sense $E_8$ lay in wait, determined by ingredients fixed earlier. The content of the theorem is therefore not \emph{that} the completion exists but that the framework's own dynamics \emph{forces} it, and forces it with exact economy: the monodromy $M$ was selected for Sol geometry and the Weinberg-angle route (Paper XX), the pillowcase by the generation structure, the quaternionic fibre by the Spin$^c$ dictionary (Paper XXI) --- none of these choices was made with $E_8$ in view --- yet the two arithmetically distinguished primes of the resulting construction, $2$ (the orbifold prime) and $\pfive$ (the golden discriminant), each supply exactly one missing generator of $2I$, with nothing left over and nothing absent, and with the $(2,3,5)$ triangle filled by three different corpus mechanisms (\S\ref{sec:235}). No continuous modulus survives into either generator.

This is rigidity evidence of the same species as the corpus's ``Two 3s'' coincidence (Paper XXII \S2.2 \texttt{thm:rigidity}; Paper XVI \S3.1; Paper XXI \S2.6): zero-parameter matches emerging from independently anchored integers. What such evidence establishes is \emph{canonicality} --- it raises the cost of dismissing $\Keight$ as an arbitrary construction. What it does not establish, and what this paper nowhere claims, is that $\Keight$ describes nature; that question is adjudicated exclusively by the corpus's numerical predictions against measurement. Derived $\neq$ confirmed.

\subsection{Not a gauge group}\label{sec:notgauge}

$E_8$ has entered fundamental physics twice as a \emph{gauge group}, with instructive fates. In the heterotic string \cite{GHMR}, $E_8 \times E_8$ is an input forced by anomaly cancellation (the rank-16 even self-dual lattices being $E_8 \oplus E_8$ and $D_{16}^+$). In the proposal of \cite{Lisi}, $E_8$ was taken as a literal unification group containing the Standard Model and gravity; this is excluded by the Distler--Garibaldi no-go theorem \cite{DistlerGaribaldi}: $E_8$ cannot accommodate three chiral fermion generations.

The present result is categorically different, and the difference is load-bearing. Here $E_8$ is an \emph{output}, not an input; it is an \emph{arithmetic lattice} (the icosian ring under the rationalised norm), not a gauge group; and the framework formally excludes $E_8$ as a fundamental gauge symmetry. In particular the Distler--Garibaldi obstruction does not apply: nothing is embedded in $E_8$ representations, and the generation count in this corpus is topological ($\chi(\Keight) = 12$ with the Spin$^c$ halving bridge \cite{SpincBridge, PaperXXI}), not representation-theoretic. The mathematical lineage of the identification used here is likewise non-stringy: the icosian ring of Conway--Sloane \cite{ConwaySloane}, the maximal-order uniqueness of Tits \cite{Tits}, the golden Galois twist of Moody--Patera \cite{MoodyPatera}, and --- one step further back --- the McKay correspondence \cite{McKay}, which pairs the binary icosahedral group $2I$ with the affine $E_8$ Dynkin diagram. What this paper adds to that lineage is a \emph{dynamical origin} for $2I$ on a physical manifold: an Anosov toral monodromy whose two modular shadows generate it.

It is also worth recording that the $E_8$ lattice itself predates all physical applications and was not constructed out of curiosity but found at the end of exhaustive classifications: as the extremal even unimodular lattice in dimension 8 in the quadratic-forms tradition (Korkine--Zolotareff \cite{KZ}; optimality of its packing proved in \cite{Viazovska}), and as the largest exceptional case of the Killing--Cartan classification. Canonicality at the receiving end is precisely what makes a forced arrival there informative.

\subsection{Remark: the golden ratio in $E_8$'s experimental signature}\label{sec:coldea}

\rtag{M, observation only} $E_8$ has to date exactly one experimental appearance in nature: the emergent $E_8$ spectrum of the critical Ising chain in a longitudinal field, predicted by Zamolodchikov \cite{Zamolodchikov} and observed in the quasi-one-dimensional ferromagnet $\mathrm{CoNb_2O_6}$ by Coldea et al.\ \cite{Coldea}. The signature ratio of its two lightest masses is
\begin{equation}
\frac{m_2}{m_1} = 2\cos\frac{\pi}{5} = \varphi = \trd(q_5).
\end{equation}
That the one number nature has so far displayed from $E_8$ is precisely the reduced trace of the canonical unit constructed in Theorem \ref{thm:q5} is reported here as an observation, under the same numerology guard as Remark \ref{rem:guard}: no mechanism is claimed, no step of this paper uses it, and it is logged as GR-6. It is recorded because both numbers have the same proximate origin --- $2\cos(\pi/5)$ as the trace of an order-10 element --- and because it identifies a concrete locus at which the outlook programme of \S\ref{sec:outlook} would meet experiment if it succeeds.

%%=====================================================================
\section{Outlook: the $E_4$ theta-series bridge}\label{sec:outlook}
%%=====================================================================

\rtag{M, programme target} The theta series of the $E_8$ root lattice is the weight-4 Eisenstein series, $\theta_{E_8}(\tau) = E_4(\tau)$; equivalently, under the Conway--Sloane rationalisation, the icosian ring carries $E_4$ as its norm generating function, with a Hilbert-modular refinement over $\Q(\sqrt5)$ available before rationalisation. The corpus's partition-function machinery on $\Keight$ (Papers XI, XV, XXII \cite{PaperXI, PaperXV, PaperXXII}) is mock modular, with completions already involving Eisenstein-type non-holomorphic corrections ($E_2^*$). The bounded derivation target is therefore: \emph{derive the appearance of $E_4$ --- or its Hilbert-modular refinement --- in the $\Keight$ partition function from the icosian identification of Theorem \ref{thm:icosian}}, so that the arithmetic lattice acquires derivational consequences in the mass/Yukawa sector rather than remaining a terminal identification. If established, the result of this paper would upgrade from rigidity evidence (\S\ref{sec:calibration}) to a load-bearing input of the empirical class; the observation of \S\ref{sec:coldea} marks where such an upgrade would first touch experiment. This is a single, well-posed target suitable for one adversarial derivation cycle; nothing in the present paper depends on it.

%%=====================================================================
\section{What this paper does and does not prove}\label{sec:doesnot}
%%=====================================================================

\subsection{Proved here}
Under the Transport Principle --- the single identified framework input, whose level-2 instance is settled corpus data --- the enlargement chain \eqref{eq:chain} is complete: the corner order is forced through the Hurwitz order (Theorem \ref{thm:hurwitz}) to the icosian ring (Theorems \ref{thm:q5}, \ref{thm:icosian}), with both forcing units constructed canonically from monodromy data, the Hypothesis H1 of the superseded Note G1 discharged (Proposition \ref{prop:h1}), and Gaps G1 and G2 of the icosian programme closed. The named analytic difficulty of the programme is concentrated in TP (GR-3) and is not hidden inside the proofs, which are elementary throughout.

\subsection{The halving pattern, partially audited}
The corpus has met \emph{halving} four times: half-period corner coordinates \eqref{eq:corners}, the Spin$^c$ lift $J^2 = -1$ \cite{SpincBridge}, the half-integer Hurwitz units (Lemma \ref{lem:rigidity}), and the coefficient $\tfrac12$ of $q_5$. The superseded Note G1 flagged at \rtag{M} the suggestion that a single Spin$^c$ mechanism underlies them all. This audit is now partially complete: within the construction, the unit-denominator instances are \emph{derived} --- the $\tfrac12$ of \eqref{eq:Uset} and of $q_5$ is in both cases the half-angle cosine of the binary double cover (Lemma \ref{lem:rigidity}; Theorem \ref{thm:q5}, Step 4), i.e.\ the same binary lift that produces $J^2 = -1$. The residual claim --- that the half-period corner lattice is itself an instantiation of the same Spin$^c$ mechanism rather than an independent geometric fact --- remains a pattern observation at \rtag{M} (GR-7); no statement above relies on it.

\subsection{Novelty delineation}
Lemmas \ref{lem:blind}--\ref{lem:rigidity}, the witnesses, and every algebraic identification are classical in substance (Hamilton; Hurwitz \cite{Hurwitz}; Skolem--Noether; the octahedral and ADE classifications; Conway--Sloane \cite{ConwaySloane}; Tits \cite{Tits}; Vign\'eras \cite{Vigneras}). No claim is made that the chain \eqref{eq:chain} is new as pure algebra, nor that any individual icosian unit is new --- all 120 are tabulated in \cite{ConwaySloane}. What is new is the \emph{forcing mechanism}: the identification of a hyperbolic toral monodromy, transported through an orbifold--quaternion dictionary by one Spin$^c$ principle, as a dynamical engine of order enlargement --- with the two enlargement steps arising as the two modular shadows of the same matrix at the two arithmetically distinguished primes, the golden axis weights sourced in the eigenbasis, the order 10 in the ramified-prime reduction, and the half-integrality in the double cover. The nearest relatives in the literature carry no dynamics: the McKay correspondence \cite{McKay} pairs $2I$ with the $E_8$ diagram combinatorially, and the quasicrystal icosians of Moody--Patera \cite{MoodyPatera} use the golden Galois twist statically. The mechanism may have applicability to other settings in which a dictionary between orbifold fixed-point data and quaternion orders is available.

%%=====================================================================
\section{Gap register}\label{sec:gaps}
%%=====================================================================

\begin{gapbox}[Gap register (honest tags; a short correct register outranks a long optimistic proof)]
\textbf{GR-1 \rtag{M} --- step normalisation (T3).} No canonical metric identification of the combinatorial 5-cycle with the geometric orbit pentagon (pentagon vs pentagram step: $2\pi/5$ vs $4\pi/5$; $q_5$ vs $q_5^{\,3}$; $\trd = \varphi$ vs $1-\varphi$). Obstruction: the transport constrains orbit structure and parity, not the cyclic embedding. Consequence bounded: $\langle q_5\rangle = \langle q_5^{\,3}\rangle$, so the generated group, the order $\Ico$, and Theorem \ref{thm:icosian} are invariant; the exhibited canonical branch has $\trd = \varphi \in \{\varphi, \varphi - 1\}$ as targeted by Corollary \ref{cor:nonmax}.

\smallskip
\textbf{GR-2 \rtag{D} --- orientation of $S^1_{\mathrm{Sol}}$.} Selecting $v_+$ presupposes the monodromy direction ($M$ vs $M^{-1}$). Reversal swaps $v_+ \leftrightarrow v_-$ and lands in the conjugate maximal order (Theorem \ref{thm:icosian}(i)); ``$\OrderMEF = \Ico$ up to conjugacy'' is invariant \rtag{R}.

\smallskip
\textbf{GR-3 \rtag{D} --- Transport Principle.} The single framework input: monodromy-equivariance of the corner dictionary with spinorial (double-cover) faithfulness. Its level-2 instance is settled (Sol Permutation Theorem $+$ $J^2 = -1$ $+$ $Q_8$ module action); its level-5 use extends the same lattice frame $\R$-linearly. Not derived here from deeper principles; the bounded derivation target (Spin$^c$ parallel transport around $S^1_{\mathrm{Sol}}$) is located in \S\ref{sec:tp}.

\smallskip
\textbf{GR-4 \rtag{G, observation only}.} $12$ pillowcase 5-torsion points $= \chi(\Keight)$. Reported per the numerology guard (Remark \ref{rem:guard}); no step uses it.

\smallskip
\textbf{GR-5 \rtag{R, remark}.} Faithful order-10 transport yields $\trd \in \{\varphi, 1-\varphi\}$; Corollary \ref{cor:nonmax}'s pair $\{\varphi$ (order 10), $\varphi - 1$ (order 5)$\}$ is harmonised by $\trd(q_5^{\,2}) = \varphi^2 - 2 = \varphi - 1$ --- the order-5 adjunction is the square of the canonical order-10 unit. No conflict.

\smallskip
\textbf{GR-6 \rtag{M, observation only}.} The Coldea--Zamolodchikov coincidence $m_2/m_1 = \varphi = \trd(q_5)$ (\S\ref{sec:coldea}). No mechanism claimed; candidate point of contact for the \S\ref{sec:outlook} programme.

\smallskip
\textbf{GR-7 \rtag{M} --- the residual halving claim.} That the half-period corner lattice is an instantiation of the same Spin$^c$ binary-lift mechanism now established for the unit denominators (\S\ref{sec:doesnot}). Pattern observation; audit deferred.

\smallskip
\textbf{Hard-constraint compliance (per the G2 brief):} mechanism delivered, not merely the specimen; no string-theoretic $E_8$ machinery (heterotic bundles, $E_8 \times E_8$ gauge sectors, lattice CFTs) imported --- $E_8$ arises solely as the Conway--Sloane arithmetic isometry of $\Ico$; Blocks A--C consumed unmodified; $\tau_2$ and all continuous moduli absent from both constructed units.
\end{gapbox}

%%=====================================================================
\section{Roadmap: results to proof sites}\label{sec:roadmap}
%%=====================================================================

\begin{center}
\small
\begin{tabular}{@{}lll@{}}
\toprule
\textbf{Result} & \textbf{Proof site (this paper)} & \textbf{External / corpus anchor} \\
\midrule
Lemma \ref{lem:ring} ($\Z[\varphi^{\pm2}] = \Zphi$) & \S\ref{sec:scalars} & Paper XX $\mathcal{M}$-7 \\
Transport Principle (TP) & \S\ref{sec:tp} (input, \rtag{D}) & Paper XXI \S4.3; Spin$^c$ Bridge \\
Lemma \ref{lem:3cycle} (Sol 3-cycle) & consumed; Remark \ref{rem:mod2} & Paper XV \S8 \texttt{thm:3cycle} \rtag{R} \\
Lemma \ref{lem:blind} (axis triviality of $Q_8$) & \S\ref{sec:level2} & Paper XXI \S4.3 ($V_{\mathrm{Aut}}$) \\
Lemma \ref{lem:rigidity} (rigidity; half-integrality) & \S\ref{sec:level2} & Skolem--Noether; $SO(3)$ classification \\
Theorem \ref{thm:hurwitz} (Hurwitz Forcing) & \S\ref{sec:level2} & --- (new; supersedes Note G1) \\
Corollary \ref{cor:nonmax} (non-maximality; hand-off) & \S\ref{sec:level2} & Conway--Sloane \S8.2; Vign\'eras \\
Witnesses W-1--W-4 & \S\ref{sec:witnesses} & elementary \rtag{R} \\
Theorem \ref{thm:q5} (canonical trace-$\varphi$ unit) & \S\ref{sec:level5} & --- (new) \\
Theorem \ref{thm:icosian} (icosian completion; arithmetic $E_8$) & \S\ref{sec:icosian} & Tits; Conway--Sloane (cite-only) \\
Proposition \ref{prop:h1} (H1 discharged) & \S\ref{sec:twoeval} & Paper III Addendum \S5; Paper XV \texttt{thm:order3} \\
$(2,3,5)$ decomposition & \S\ref{sec:235} & Papers XV, XXI $+$ this paper \\
Coldea remark \rtag{M} & \S\ref{sec:coldea} & Zamolodchikov; Coldea et al. \\
$E_4$ outlook \rtag{M} & \S\ref{sec:outlook} & Papers XI, XV, XXII \\
\bottomrule
\end{tabular}
\end{center}

\section*{Acknowledgements}
The author used Claude (Anthropic) and Gemini (Google DeepMind) as mathematical development and analytical tools during the preparation of this work. All intellectual content, theoretical framework, and scientific claims are the sole responsibility of the author.

%%=====================================================================
%% References
%%=====================================================================
\begin{thebibliography}{99}

\bibitem{PaperXV} D.~Master, \emph{The Geometric Origin of Mock Modularity}, Master Equation Framework series (2026). \S8: Theorem \texttt{thm:3cycle} (Sol Permutation Theorem); Theorem \texttt{thm:order3}; \S8.5 Galois--Orbifold Correspondence.

\bibitem{PaperXX} D.~Master, \emph{New Axioms for Old Problems}, Master Equation Framework series (2026). $\Sigma$-2g, $\mathcal{M}$-7: $M = \left(\begin{smallmatrix}2&1\\1&1\end{smallmatrix}\right)$, eigenvalues $\varphi^{\pm2}$, Weyl discriminant $5 = \disc\Q(\sqrt5)$.

\bibitem{PaperXXI} D.~Master, \emph{The Quaternionic Unification on $\Keight$}, Master Equation Framework series (2026). \S4.1--4.3: corners, dictionary, $Q_8$ module action, $V_{\mathrm{Aut}}$; \S4.4: Spin$^c$ as $Sp(1) \to U(1)$ reduction; \S2.6: $\lambda_{0,2} = \pi$ at $c_1(L) = 3$, $\tau_2$-independent.

\bibitem{PaperIIIAdd} D.~Master, \emph{Paper III Addendum: Sol Monodromy}, Master Equation Framework series (2026). \S5: corner architecture and coordinate assignment ($C_2 \leftrightarrow \C_i$, $C_1 \leftrightarrow \C_j$, $C_3 \leftrightarrow \C_k$).

\bibitem{PaperXI} D.~Master, \emph{The Mock Modular Partition Function of $\Keight$}, Master Equation Framework series (2026).

\bibitem{PaperXXII} D.~Master, \emph{The Almost-Quaternionic / Mock Modular Correspondence}, Master Equation Framework series (2026). \S2.2 \texttt{thm:rigidity} (``Two 3s'' rigidity).

\bibitem{SpincBridge} D.~Master, \emph{Spin$^c$ Generation Halving Bridge}, Master Equation Framework corpus working note (2026). $J^2 = -1$.

\bibitem{ConwaySloane} J.~H.~Conway and N.~J.~A.~Sloane, \emph{Sphere Packings, Lattices and Groups}, 3rd ed., Springer (1999), ch.~8 \S2.1 (icosians; Hurwitz units; rationalised norm and the $E_8$ identification).

\bibitem{Vigneras} M.-F.~Vign\'eras, \emph{Arithm\'etique des alg\`ebres de quaternions}, Lecture Notes in Mathematics 800, Springer (1980) (orders, discriminants, maximality, Skolem--Noether).

\bibitem{Tits} J.~Tits, Quaternions over $\Q(\sqrt5)$, Leech's lattice and the sporadic group of Hall--Janko, \emph{J.\ Algebra} \textbf{63} (1980) 56--75 (the icosian ring as the unique maximal order; $E_8$ identification).

\bibitem{MoodyPatera} R.~V.~Moody and J.~Patera, Quasicrystals and icosians, \emph{J.\ Phys.\ A} \textbf{26} (1993) 2829--2853 (golden Galois twist; $E_8 \to H_4 \oplus \varphi H_4$ shadow split).

\bibitem{Hurwitz} A.~Hurwitz, \emph{Vorlesungen \"uber die Zahlentheorie der Quaternionen}, Springer (1919).

\bibitem{McKay} J.~McKay, Graphs, singularities, and finite groups, in \emph{Proc.\ Symp.\ Pure Math.}\ \textbf{37}, AMS (1980) 183--186 (the correspondence $2I \leftrightarrow$ affine $E_8$).

\bibitem{KZ} A.~Korkine and G.~Zolotareff, Sur les formes quadratiques, \emph{Math.\ Ann.}\ \textbf{6} (1873) 366--389 (the $E_8$ lattice as extremal form).

\bibitem{Viazovska} M.~Viazovska, The sphere packing problem in dimension 8, \emph{Ann.\ of Math.}\ \textbf{185} (2017) 991--1015.

\bibitem{GHMR} D.~J.~Gross, J.~A.~Harvey, E.~Martinec and R.~Rohm, Heterotic string, \emph{Phys.\ Rev.\ Lett.}\ \textbf{54} (1985) 502--505.

\bibitem{Lisi} A.~G.~Lisi, An exceptionally simple theory of everything, arXiv:0711.0770 (2007).

\bibitem{DistlerGaribaldi} J.~Distler and S.~Garibaldi, There is no ``Theory of Everything'' inside $E_8$, \emph{Comm.\ Math.\ Phys.}\ \textbf{298} (2010) 419--436.

\bibitem{Zamolodchikov} A.~B.~Zamolodchikov, Integrals of motion and $S$-matrix of the (scaled) $T = T_c$ Ising model with magnetic field, \emph{Int.\ J.\ Mod.\ Phys.\ A} \textbf{4} (1989) 4235--4248 (the $E_8$ mass spectrum; $m_2/m_1 = 2\cos(\pi/5)$).

\bibitem{Coldea} R.~Coldea et al., Quantum criticality in an Ising chain: experimental evidence for emergent $E_8$ symmetry, \emph{Science} \textbf{327} (2010) 177--180.

\end{thebibliography}

%%=====================================================================
\appendix
\section{BF-register audit (18 devices)}\label{app:bf}
%%=====================================================================

\begin{center}
\small
\begin{tabular}{@{}clll@{}}
\toprule
\# & \textbf{Device} & \textbf{Verdict} & \textbf{Site} \\
\midrule
1 & Motivate-before-define & Present & \S\ref{sec:intro} \\
2 & Type-first chained definitions & Present & \S\ref{sec:setting} \\
3 & Immediate gloss of sub-components & Present & \S\ref{sec:geometry}, glosses throughout \\
4 & Simplifying assumption flagged $+$ lifting located & Present & TP, \S\ref{sec:tp}; GR-3 \\
5 & Why-this-hypothesis-matters & Present & \S\ref{sec:tp} \\
6 & Numbered displays, domain$\to$codomain prose & Present & \eqref{eq:dictionary}, \eqref{eq:frame} \\
7 & Hypotheses front-loaded into statements & Present & Theorems \ref{thm:hurwitz}, \ref{thm:q5} open ``Under TP'' \\
8 & Every statement-symbol pre-defined & Present & \S\ref{sec:setting} $+$ \eqref{eq:varpi} \\
9 & Restriction regime named in statement & Present & norm-one groups, $\Zphi$ regime, \S\ref{sec:orders} \\
10 & One-line crux $+$ machinery & Present & \S\ref{sec:intro} \\
11 & Prose interpretation of bare identities & Present & glosses of \eqref{eq:Uset}--\eqref{eq:lifts}, Theorem \ref{thm:hurwitz} \\
12 & Forward reference to precise version / proof site & Present & Corollary \ref{cor:nonmax} hand-off; \S\ref{sec:outlook} \\
13 & Analytic difficulty named, not hidden & Present & \S\ref{sec:tp} status; GR-3 \\
14 & Heuristic flagged then audited & Partial & halving: derived for unit denominators (\S\ref{sec:doesnot}); GR-7 residual \\
15 & Parallel structure made explicit & Present & \eqref{eq:chain}; Remark \ref{rem:mod2} $\parallel$ W-1; \S\ref{sec:twoeval} \\
16 & Novelty delineated against prior work & Present & \S\ref{sec:doesnot}; \S\ref{sec:notgauge} (McKay, Moody--Patera, Lisi, heterotic) \\
17 & Difference-from-previous-case flagged & Present & parity visible at $\pfive$ vs invisible at $2$, \S\ref{sec:twoeval} \\
18 & Roadmap mapping results to proof sites & Present & \S\ref{sec:roadmap} \\
\bottomrule
\end{tabular}
\end{center}

\noindent\textbf{Counts: Present 17 \,$\cdot$\, Partial 1 \,$\cdot$\, Not applicable 0.}

\smallskip
\noindent The single Partial (device 14) is deliberate: the halving heuristic's audit is now half-closed --- the unit-denominator instances are derived as the double-cover half-angle (Lemma \ref{lem:rigidity}; Theorem \ref{thm:q5} Step 4) --- while the residual identification of the half-period corner lattice with the same Spin$^c$ mechanism is GR-7 and belongs to the TP derivation cycle. No numbered statement depends on the deferred item.

\bigskip
\begin{center}
\small\itshape Prepared as a standalone derivation paper within the Master Equation Framework corpus.\\
PI: Dhiren J.~Master \,$\cdot$\, June 2026 \,$\cdot$\, UK English throughout.
\end{center}

\end{document}
